\documentclass{article}
\usepackage[utf8]{inputenc}
\usepackage{amsmath, amssymb, amsthm}
\usepackage{tikz-cd}
\usepackage{hyperref}

\theoremstyle{definition}
\newtheorem{definition}{Definition}[section]
\newtheorem{example}{Example}[section]
\newtheorem{remark}{Remark}[section]

\title{A Pedagogical Introduction to Equalizers, Kernels, Limits, and Colimits}
\author{Compiled from Category Theory Notes}
\date{}

\begin{document}

\maketitle

\tableofcontents

\section{Motivation: Solving Equations in Sets}
In set theory, we often encounter the problem of finding elements where two functions agree. Given $f, g : A \to B$, the \emph{solution set} is
\[
\{ x \in A \mid f(x) = g(x) \}.
\]
This subset comes with an inclusion map $i : \{ x \in A \mid f(x) = g(x) \} \hookrightarrow A$, and it has a special property: any function $h : Z \to A$ that satisfies $f \circ h = g \circ h$ must factor uniquely through $i$. Category theory abstracts this idea into the concept of an \textbf{equalizer}.

\section{Equalizers}
\subsection{Definition via Universal Property}
Let $\mathcal{C}$ be a category and let $f, g : X \to Y$ be two parallel morphisms.

\begin{definition}
An \textbf{equalizer} of $f$ and $g$ is a pair $(E, \mathrm{eq})$ consisting of an object $E$ and a morphism $\mathrm{eq} : E \to X$ such that:
\begin{enumerate}
    \item $f \circ \mathrm{eq} = g \circ \mathrm{eq}$,
    \item For any object $Z$ and morphism $m : Z \to X$ with $f \circ m = g \circ m$, there exists a \emph{unique} morphism $u : Z \to E$ making the following diagram commute:
    \[
    \begin{tikzcd}
    Z \arrow[dr, "m"] \arrow[d, dashed, "\exists!\, u"] & \\
    E \arrow[r, "\mathrm{eq}"'] & X \arrow[r, shift left, "f"] \arrow[r, shift right, "g"'] & Y
    \end{tikzcd}
    \]
\end{enumerate}
\end{definition}

The morphism $\mathrm{eq}$ is necessarily a \textbf{monomorphism} (i.e. satisfying left elimination, it is injective in many concrete categories). The equalizer is unique up to unique isomorphism.

\subsection{Examples}
\begin{example}[In $\mathbf{Set}$]
For sets $X, Y$ and functions $f, g : X \to Y$, the equalizer object is
\[
E = \{ x \in X \mid f(x) = g(x) \},
\]
and $\mathrm{eq} : E \hookrightarrow X$ is the subset inclusion. The universal property says: any function $m : Z \to X$ whose image lies in $E$ factors uniquely through $E$.
\end{example}

\begin{example}[In $\mathbf{Grp}$]
Given group homomorphisms $f, g : G \to H$, the equalizer is the subgroup
\[
E = \{ x \in G \mid f(x) = g(x) \},
\]
with inclusion homomorphism. This is precisely the kernel of the difference map $f - g$ when $H$ is abelian.
\end{example}

\section{Kernels as Special Equalizers}
In categories equipped with a \textbf{zero object} (and thus zero morphisms $0_{X,Y} : X \to Y$), we can specialize the equalizer construction.

\begin{definition}
Let $\mathcal{C}$ be a category with zero morphisms. The \textbf{kernel} of a morphism $f : X \to Y$ is the equalizer of $f$ and the zero morphism $0_{X,Y} : X \to Y$:
\[
\mathrm{Ker}(f) = \mathrm{Eq}(f, 0_{X,Y}).
\]
\end{definition}

Diagrammatically:
\[
\begin{tikzcd}
\mathrm{Ker}(f) \arrow[r, "\mathrm{ker}"] & X \arrow[r, shift left, "f"] \arrow[r, shift right, "0"'] & Y
\end{tikzcd}
\]
Universality of the kernel:
\[
\begin{tikzcd}
Z \arrow[dr, "m"] \arrow[d, dashed, "\exists!\, u"] & \\
\mathrm{Ker}(f) \arrow[r, "\mathrm{ker}"'] & X \arrow[r, shift left, "f"] \arrow[r, shift right, "0"'] & Y
\end{tikzcd}
\]

\begin{example}[In $\mathbf{Vect}_k$]
For a linear map $T : V \to W$, the kernel is the subspace $\{ v \in V \mid T(v) = 0 \}$. The inclusion map is the equalizer.
\end{example}

\section{Coequalizers and Cokernels}
By \textbf{reversing all arrows}, we obtain the dual notions.

\subsection{Coequalizer}
\begin{definition}
A \textbf{coequalizer} of $f, g : X \to Y$ is a pair $(Q, q)$ where $q : Y \to Q$ satisfies $q \circ f = q \circ g$, and for any $(Z, h : Y \to Z)$ with $h \circ f = h \circ g$, there exists a unique $v : Q \to Z$ making the diagram commute:
\[
\begin{tikzcd}
X \arrow[r, shift left, "f"] \arrow[r, shift right, "g"'] & Y \arrow[r, "q"] \arrow[dr, "h"'] & Q \arrow[d, dashed, "\exists!\, v"] \\
& & Z
\end{tikzcd}
\]
\end{definition}

In $\mathbf{Set}$, the coequalizer of $f, g : X \to Y$ is the quotient set $Y/{\sim}$, where $\sim$ is the equivalence relation generated by $f(x) \sim g(x)$ for all $x \in X$.

\subsection{Cokernel}
\begin{definition}
In a category with zero morphisms, the \textbf{cokernel} of $f : X \to Y$ is the coequalizer of $f$ and $0_{X,Y}$:
\[
\mathrm{Coker}(f) = \mathrm{Coeq}(f, 0_{X,Y}).
\]
\end{definition}

Diagram:
\[
\begin{tikzcd}
X \arrow[r, shift left, "f"] \arrow[r, shift right, "0"'] & Y \arrow[r, "\mathrm{coker}"] & \mathrm{Coker}(f)
\end{tikzcd}
\]
Universality:
\[
\begin{tikzcd}
X \arrow[r, shift left, "f"] \arrow[r, shift right, "0"'] & Y \arrow[r, "\mathrm{coker}"] \arrow[dr, "h"'] & \mathrm{Coker}(f) \arrow[d, dashed, "\exists!\, v"] \\
& & Z
\end{tikzcd}
\]

\begin{example}[In $\mathbf{Vect}_k$]
The cokernel of a linear map $T : V \to W$ is the quotient space $W / \operatorname{Im}(T)$. The projection map is the coequalizer.
\end{example}

\section{The Big Picture: Limits and Colimits}
Both equalizers and kernels are instances of a far-reaching concept: \textbf{limits}. Dually, coequalizers and cokernels are \textbf{colimits}.

\subsection{Cones and Limits}
Let $J$ be a small ``index'' category (think of it as a pattern), and let $D : J \to \mathcal{C}$ be a \textbf{diagram} of shape $J$ in $\mathcal{C}$.

\begin{definition}
A \textbf{cone} over $D$ consists of an object $L$ (the apex) and a family of morphisms $\pi_X : L \to D(X)$ for every object $X$ in $J$, such that for every morphism $u : X \to Y$ in $J$, the following triangle commutes:
\[
\begin{tikzcd}
& L \arrow[dl, "\pi_X"'] \arrow[dr, "\pi_Y"] & \\
D(X) \arrow[rr, "D(u)"'] & & D(Y)
\end{tikzcd}
\]
\end{definition}

A \textbf{limit} of $D$, denoted $\lim D$, is a \textbf{universal cone}: a cone $(L, \pi)$ such that for any other cone $(Z, \tau)$ over $D$, there exists a \emph{unique} morphism $u : Z \to L$ with $\tau_X = \pi_X \circ u$ for all $X$.
\[
\begin{tikzcd}
Z \arrow[drr, bend left, "\tau_Y"] \arrow[ddr, bend right, "\tau_X"'] \arrow[dr, dashed, "\exists!\, u"] & & \\
& \lim D \arrow[r, "\pi_Y"] \arrow[d, "\pi_X"'] & D(Y) \\
& D(X) \arrow[ur, "D(u)"'] &
\end{tikzcd}
\]

\begin{example}[Equalizer as a Limit]
The equalizer of $f, g : X \to Y$ is the limit of the diagram
\[
\begin{tikzcd}
X \arrow[r, shift left, "f"] \arrow[r, shift right, "g"'] & Y
\end{tikzcd}
\]
The index category $J$ has two objects and two parallel non-identity arrows. The limit cone picks out the subobject $E$ with its map to $X$ (other trivial arrows are usually not drawn).
\end{example}

\begin{example}[Product as a Limit]
If $J$ is a discrete category with two objects and no non-identity arrows, a diagram $D : J \to \mathcal{C}$ is just a pair of objects $A, B$. The limit is the product $A \times B$.
\end{example}

\subsection{Cocones and Colimits}
Dualizing everything yields colimits.

\begin{definition}
A \textbf{cocone} under $D$ consists of an object $C$ and morphisms $\iota_X : D(X) \to C$ such that for every $u : X \to Y$ in $J$, the following commutes:
\[
\begin{tikzcd}
D(X) \arrow[rr, "D(u)"] \arrow[dr, "\iota_X"'] & & D(Y) \arrow[dl, "\iota_Y"] \\
& C &
\end{tikzcd}
\]
\end{definition}

The \textbf{colimit} of $D$, denoted $\operatorname{colim} D$, is a \textbf{universal cocone}: a cocone $(C, \iota)$ such that for any other cocone $(Z, \tau)$, there is a unique $v : C \to Z$ with $\tau_X = v \circ \iota_X$.
\[
\begin{tikzcd}
D(X) \arrow[rr, "D(u)"] \arrow[dr, "\iota_X"'] \arrow[ddr, bend right, "\tau_X"'] & & D(Y) \arrow[dl, "\iota_Y"] \arrow[ddl, bend left, "\tau_Y"] \\
& \operatorname{colim} D \arrow[d, dashed, "\exists!\, v"] & \\
& Z &
\end{tikzcd}
\]

\begin{example}[Coequalizer as a Colimit]
The coequalizer of $f, g : X \to Y$ is the colimit of the parallel-arrow diagram. The universal cocone gives the quotient object.
\end{example}

\begin{example}[Coproduct as a Colimit]
For a discrete diagram $\{A, B\}$, the colimit is the coproduct $A \sqcup B$.
\end{example}

\section{Summary Table}
\begin{center}
\begin{tabular}{c|c|c}
\textbf{Construction} & \textbf{Limit or Colimit?} & \textbf{Diagram Shape} \\
\hline
Equalizer & Limit & $\bullet \rightrightarrows \bullet$ \\
Kernel & Limit & $\bullet \rightrightarrows \bullet$ (one arrow zero) \\
Product & Limit & $\bullet \quad \bullet$ (discrete) \\
Pullback & Limit & $\bullet \to \bullet \leftarrow \bullet$ \\
Coequalizer & Colimit & $\bullet \rightrightarrows \bullet$ \\
Cokernel & Colimit & $\bullet \rightrightarrows \bullet$ (one arrow zero) \\
Coproduct & Colimit & $\bullet \quad \bullet$ (discrete) \\
Pushout & Colimit & $\bullet \leftarrow \bullet \to \bullet$
\end{tabular}
\end{center}

\section*{Final Remarks}
The language of limits and colimits unifies a vast array of constructions across algebra, topology, and logic. Mastering equalizers and coequalizers is the first step toward understanding this elegant framework.


\end{document}